The real count is Fact \claimref{tower.witt-index-quantum-capacity}. For the complex count (the modes of the Witt reading of \claimref{one-relation.two-presentations}): a totally isotropic $U \subset V_{\C}$ satisfies $U \subseteq U^{\perp}$, and nondegeneracy gives $\dim U + \dim U^{\perp} = d$, so $2\dim U \leq d$; the bound is attained by the null spans of $\lfloor d/2 \rfloor$ pairs, built from the real hyperbolic planes together with pairs $(e + I_{\hat V}e')/2$, $-(e - I_{\hat V}e')/2$ on each two-dimensional block of the anisotropic kernel, where $V_{\C} = V \oplus I_{\hat V}V$ is realised in the odd extension $\hat V$ of $V$ (of $V \oplus \R f$, for one further vector $f$, when $d$ is odd), whose central pseudoscalar $I_{\hat V}$ squares to $-1$ and commutes with every vector (Theorem \claimref{std-algebra.complexification-is-the-odd-extension}); their square $\{c(e+I_{\hat V}e'), c(e+I_{\hat V}e')\} = 2(Q(e) - Q(e')) = 0$ (\claimref{std-algebra.vector-squares}) vanishes because the kernel is definite and its diagonal entries equal. Both counts are exhibited by the script on $(1,1)$, $(1,3)$, $(2,4)$: the pairs are written down and every anticommutator computed. (1) $\min(1,3) = 1$; the anisotropic complement of the boost plane is the negative-definite $(e_{2}, e_{3})$ block of the Witt decomposition (\claimref{std-algebra.witt-decomposition}), on which the script confirms $Q(ae_{2} + be_{3}) = -(a^{2} + b^{2}) < 0$ for rational $(a,b) \neq 0$, so no null vector lies in the anisotropic complement and no second anticommuting real pair exists; the Witt decomposition is unique up to isometry by Witt's extension theorem (Fact \claimref{one-relation.imported-witts-extension-theorem}), which is the sense in which the boost plane is unique, and the same Witt fact is the uniqueness of the invariant ratio (Lemma \claimref{one-relation.third-constant-conversion-constant-between-charts}(ii)). The agreement locus: for $p \leq q$, $r_{\R} = p$ and $r_{\C} = \lfloor d/2 \rfloor = \lfloor (p+q)/2 \rfloor$, so $r_{\R} = r_{\C}$ if and only if $2p \leq p + q \leq 2p + 1$, that is $0 \leq q - p \leq 1$, and symmetrically for $q \leq p$; the script's grid check illustrates this. (2) $\min(p+1,q+1) = \min(p,q) + 1$ and $\lfloor (d+2)/2 \rfloor = \lfloor d/2 \rfloor + 1$, arithmetic identities, recorded by the script on the grid $p, q \leq 8$. (3) $\lfloor 6/2 \rfloor = 3$ and $\min(2,4) = 2$, with the three pairs of $\Cl(2,4)_{\C}$ exhibited, two real and one complex.
